\hypertarget{r29-ao1-ao1-forward-hnm-finite-wilson-second-moment-bound}{%
\section{HNM AO1: finite Wilson second-moment bound\workstar}\label{r29:ao1:r29-ao1-ao1-forward-hnm-finite-wilson-second-moment-bound}}
% BEGIN CURRENT REGISTRY NOTE
\paragraph{Statement alias HNM-P-AO1.}\label{stmt:AO1} This identifier denotes the scoped mathematical claim admitted by this chapter\textquotesingle s gate. It adds no assumption and establishes no literature-priority claim.
\begin{quote}\small\textbf{Current extension: HNM-C-AO2.} The historical limitation below remains the scope of this original result. The later, separately gated statement in Section~\ref{stmt:AO2} records the extension: For the same actual I1/AJ1/AK2 orthant physical state, chi belongs to D(H\_phys), its second spectral moment is at most alpha\textasciicircum{}2(36+28|tau|)<37alpha\textasciicircum{}2, and its first moment equals alpha omega(1-W\textasciicircum{}2). Uniform first-energy tails bounded by B2/R close the former AK2 moment-loss limitation. The proof uses the actual centered spectral measures, compact spectral tests and bounded-local expectations. Its own model and state assumptions remain required.\end{quote}
% END CURRENT REGISTRY NOTE

\begin{quote}\small \textbf{Admission: accepted within scope.} For the same actual origin xz Wilson and every finite whole-star orthant volume containing its complete region, the centered physical second moment obeys \(\mu_2\le \alpha^{2}(36+28|\tau|)\le \alpha^{2}(36+7/16384)<37\alpha^{2}\). The proof derives the kinetic bound from the actual selected-strip reference, retains the exact ground scalar before its legitimate upper-bound relaxation, and proves the finite domain-valid electric commutator including its derivative cross term.\par The typeset body below preserves the source-bound forward derivation. Its prospective language is historical; this boxed gate records the final reviewed verdict. The separately authored reverse and skeptical records remain linked. The common admission and attributed refinements are determined by the gate, not by a renamed equation.\end{quote}
Human project author: \textbf{Hruday N M (BUNZEEY)}. Independent forward derivation of the shared prospective commutator route, without reading current reverse or skeptic results. The method uses standard product rules; scientific priority is unverified.

\textbf{Finite-volume result:} the original origin xz Wilson in every contracted actual whole-star orthant ground state satisfies

\begin{equation}
\int E^2d\nu_\Lambda(E)=\|K_\Lambda\chi_\Lambda\|^2
\le\alpha^2(36+28|\tau|)
\le\alpha^2(36+7/16384)<37\alpha^2.\tag{HNM-AO1.1}
\label{r29:ao1:eq1}
\end{equation}

Here \(\chi_\Lambda=(W-\omega_\Lambda(W))\Omega_\Lambda\) and \(K_\Lambda=H^{\rm raw}_\Lambda-E^{\rm raw}_\Lambda\) is the actual centered physical operator. The inherited \(|\tau|<\tau_*\) and \(|\tau|\le 2^{-16}\) remain in force. Reference-ground energy is not an excited-vector spectral moment.

\hypertarget{r29-ao1-exact-local-kinetic-reconstruction}{%
\subsection{Exact local kinetic reconstruction}\label{r29:ao1:r29-ao1-exact-local-kinetic-reconstruction}}

The complete cover is \(R=\{0,e_z\}\): 48 links and 36 endpoint actions. Its actual orthant incident anchor union is \((R-S)\cap\mathbb Z_+^3=\{0,e_z\}\), counting the star at zero once. AK1 therefore gives \(\omega_\Lambda(h_R)\le 28|\tau|\).

Let \(T_R=\alpha\sum_{e\text{ owned by }R}C_e\), \(V_R^{\rm sel}=-\sum_{\substack{f\text{ selected}\\ f\text{ in }R}}\lambda_fW_f\), and \(B_R=\sum_{b\in R}E_{{\rm strip},b}\). Keeping every selected ground scalar,

\begin{equation}
\delta h_R=T_R+V_R^{\mathrm{sel}}-B_R,\qquad
T_R=\delta h_R-V_R^{\mathrm{sel}}+B_R,\qquad\delta=\alpha/8.\tag{HNM-AO1.2}
\label{r29:ao1:eq2}
\end{equation}

The constant Haar trial on each strip has zero electric and selected-Wilson mean, so \(E_{\mathrm{strip},b}\le 0\). Each strip has selected absolute budget at most \(\alpha/2+\alpha/8+\alpha/2=9\alpha/8\). Thus

\begin{equation}
\omega_\Lambda(T_R)\le{\alpha\over8}28|\tau|+{9\alpha\over4}
=\alpha(9/4+(7/2)|\tau|).\tag{HNM-AO1.3}
\label{r29:ao1:eq3}
\end{equation}

Both coupling signs and every allowed signed selected coefficient are covered. Omitting \(B_R\) from the exact identity is wrong; discarding its nonpositive contribution after writing an upper inequality is legitimate. Omitting \(V_{R}^{\mathrm{sel}}\) from the reconstruction is also wrong. At \(\tau=0\) a nontrivial selected ground has zero reference energy while its electric expectation need not vanish.

\hypertarget{r29-ao1-domains-and-four-original-links}{%
\subsection{Domains and four original links}\label{r29:ao1:r29-ao1-domains-and-four-original-links}}

At finite volume the Casimir sum is elliptic on compact \(SU(2)^L\), and the full magnetic potential is a bounded smooth real multiplier. Its operator domain is the product H² domain. The ground is smooth. W has bounded derivatives through order two, so multiplication preserves this domain. Hence \(\chi_\Lambda\) lies in \(D(K_\Lambda)\), and its second spectral moment equals the squared generator norm. No infinite-volume domain is assumed.

All selected and omitted magnetic multipliers commute with W. For \(C_e=-\sum_aX_{ea}^2\),

\([C_e,W]u=(C_eW)u-2\sum_a(X_{ea}W)X_{ea}u\).

The original four stored links are distinct. Each contributes \(C_eW=3W/4\), with inverse occurrences retaining their derivative side and sign. Consequently

\begin{equation}
K_\Lambda\chi_\Lambda=3\alpha W\Omega_\Lambda-2\alpha Z_\Lambda,
\quad Z_\Lambda=\sum_{e\in p,a}(X_{ea}W)X_{ea}\Omega_\Lambda.\tag{HNM-AO1.4}
\label{r29:ao1:eq4}
\end{equation}

AK2's half-Pauli identity gives \(\sum_{e\in p,a}(X_{ea}W)^2=1-W^2\le1\). Pointwise Cauchy--Schwarz followed by integration yields

\begin{equation}
\|Z_\Lambda\|^2\le\int(1-W^2)\sum_{e\in p,a}|X_{ea}\Omega_\Lambda|^2
\le\omega_\Lambda(T_R)/\alpha.\tag{HNM-AO1.5}
\label{r29:ao1:eq5}
\end{equation}

The derivative cross term is retained: the actual selected reference is generally not constant. Using \(\|W\Omega\|\le1\) and \(\|u+v\|^2\le2\|u\|^2+2\|v\|^2\),

\begin{equation}
\|K_\Lambda\chi_\Lambda\|^2\le18\alpha^2+8\alpha\omega_\Lambda(T_R)
\le\alpha^2(36+28|\tau|).\tag{HNM-AO1.6}
\label{r29:ao1:eq6}
\end{equation}

This proves the claimed bound in physical energy-squared units. A frequency second moment would divide by hbar². The fixed energy reference and clock are unchanged.

\hypertarget{r29-ao1-controls-and-finite-scope}{%
\subsection{Controls and finite scope}\label{r29:ao1:r29-ao1-controls-and-finite-scope}}

At the fully free corner, where all selected coefficients and \(\tau\) vanish, the Haar ground is constant, variance is 1/4, and the Wilson energy is 3alpha. Its exact second moment is 9alpha²/4 despite zero reference-ground energy. This rejects equating those two energies.

The checker reconstructs origin versus interior incidence (two versus seven stars), the complete link cover, and exact coefficients. A smooth one-rotor diagnostic \(\Omega(q)=1+q_0/3\), \(W=q_0\) at \(q_0=3/5\) has a nonzero gradient cross term; dropping it changes the commutator. It is an algebra control, not an I1 ground. A finite matrix expectation fixture tests the negative ground scalar in the exact reconstruction and its safe removal only for an upper inequality. Doubling Lie generators quadruples squared gradients and is rejected. Measures with bounded first moments but escaping high-energy tails have growing second moments, so first-moment boundedness cannot replace this proof.

The regional constant is specific to the original orthant cover; interior or summable-model substitution needs a new dictionary. No limit is taken. Uniform integrability, limiting moment equality and the actual infinite-volume operator domain remain for adaptive AO2.

Reproduce from the repository root:
\begin{verbatim}
python -B research/round29/forward/ao1/check.py \
  --output /absolute/new/output
\end{verbatim}

\subsection*{Final admitted limits and evidence}
\begin{itemize}
\item This is a finite-volume uniform moment theorem only; infinite operator-domain membership and first-moment equality require a separate limit argument.
\item The actual orthant state and both inherited conditions \(|\tau|<\tau_*\) and \(|\tau|\le 2^{-16}\) are retained; this does not evaluate \(\tau_*\).
\item The two-star origin incidence cannot be substituted for an interior region or a different summable model.
\item Bounded second moments do not imply equality of limiting second moments. Standard commutator methods and the shared prospective panel route are acknowledged; scientific priority remains unverified.
\item Fixed positive physical scales and original half-Pauli normalization are essential. No continuum or experimental claim follows.
\end{itemize}
\repo{research/round29/advisor/ao1-gate.json}\par\repo{research/round29/reverse/ao1/report.md}\par\repo{research/round29/skeptic/ao1.md}

\emph{Sequencing note.} Prospective statements in this frozen derivation describe the moment of its admission. The final Round29 roadmap below governs subsequent work.
